\subsection{约化范数与约化迹的传递性}

命题 7. —— 设 L 是 A 的一个极大交换半单子代数，a 是 L 的一个元素。我们有

\[
\mathrm{Pcrd} _ {\mathrm{A/K}} (a; \mathrm{X}) = \mathrm{Pc} _ {\mathrm{L/K}} (a; \mathrm{X}),\tag{52}
\]

\[
\mathrm{Trd} _ {\mathrm{A/K}} (a) = \mathrm{Tr} _ {\mathrm{L/K}} (a),\tag{53}
\]

\[
\mathrm{Nrd} _ {\mathrm{A/K}} (a) = \mathrm{N} _ {\mathrm{L/K}} (a).\tag{54}
\]

由 VIII, 第 262 页，命题 3，L-模 A 与 \(L^{n}\) 同构；因此我们有关系式

\[
\mathrm{Pc} _ {\mathrm{A/K}} (a; \mathrm{X}) = \mathrm{Pc} _ {\mathrm{L/K}} (a; \mathrm{X}) ^ {n}.
\]

由于多项式 \(\mathrm{Pc}_{\mathrm{L}/\mathrm{K}}(a;\mathrm{X})\) 是首一的，它于是等于约化特征多项式 \(\mathrm{Pcrd}_{\mathrm{A}/\mathrm{K}}(a;\mathrm{X})\)（VIII, 第 339 页，引理 2）；这就给出公式 (52)。比较 (52) 两边 \(X^{n-1}\) 的系数（相应地，常数项），我们便得到公式 (53)（相应地，(54)）。

推论. —— 设 D 是 K 上的一个有限次数的体，其中心为 K。设 a 是 K 的一个元素，L 是 D 的一个包含 a 的极大交换子域。我们有

\[
\begin{array}{c} \mathrm {Pcrd_ {D / K}} (a; \mathrm{X}) = \mathrm {Pc_ {L / K}} (a; \mathrm{X}), \\ \mathrm{Tr} _ {\mathrm{D/K}} (a) = \mathrm{Tr} _ {\mathrm{L/K}} (a), \\ \mathrm{Nrd} _ {\mathrm{D/K}} (a) = \mathrm{N} _ {\mathrm{L/K}} (a). \end{array}\tag{55}
\]

事实上，由 VIII, 第 265 页，推论 2，D 的极大交换子域 L 是 D 的一个极大交换半单子代数。

命题 8. —— 设 B 是 A 的一个单子代数。用 L 表示 B 的中心，用 B' 表示 B 在 A 中的交换子。那么 B' 是体 L 上的一个中心单代数；我们用 r 表示它的约化次数。对 B 的每个元素 b，我们有关系式

\[
\mathrm{Pcrd} _ {\mathrm{A/K}} (b; \mathrm{X}) = \mathrm{N} _ {\mathrm{L[X]/K[X]}} (\mathrm{Pcrd} _ {\mathrm{B/L}} (b; \mathrm{X})) ^ {r},\tag{56}
\]

\[
\mathrm{Trd} _ {\mathrm{A/K}} (b) = r \mathrm{Tr} _ {\mathrm{L/K}} (\mathrm{Trd} _ {\mathrm{B/K}} (b)),\tag{57}
\]

\[
\mathrm{Nrd} _ {\mathrm{A/K}} (b) = \mathrm{N} _ {\mathrm{L/K}} (\mathrm{Nrd} _ {\mathrm{B/L}} (b)) ^ {r}.\tag{58}
\]

引理 6. —— 设 \(\mathbf{K}'\) 是一个有限次数为 \(d\) 的交换 \(\mathbf{K}\)-代数，并设

\[
\mathrm{P} (\mathrm{X}) = \mathrm{X} ^ {s} + a _ {1} \mathrm{X} ^ {s - 1} + \dots + a _ {s}
\]

是一个系数属于 \(K'\) 的首一多项式。那么多项式 \(Q = N_{K'[X]/K[X]}(P)\) 在 \(K[X]\) 中是首一的，次数为 sd，\(X^{sd-1}\) 在 \(Q(X)\) 中的系数等于 \(\mathrm{Tr}_{\mathrm{K}'/\mathrm{K}}(a_1)\)，且 Q 的常数项为 \(N_{K'/K}(a_s)\)。

我们把 \(K^{\prime}\)-代数 \(\mathrm{K}^{\prime}[\mathrm{T}]/( \mathrm{P}(\mathrm{T}))\) 记作 \(K^{\prime\prime}\)，并把 T 在 \(K^{\prime\prime}\) 中的典范类记作 t。序列 \((1, t, \ldots, t^{s-1})\) 是 \(K^{\prime\prime}\) 在 \(K^{\prime}\) 上的一组基，而由 t 作乘法在这组基下的矩阵形如

\[
\tau = \left( \begin{array}{c c c c c} 0 & 0 & \dots & 0 & - a _ {s} \\ 1 & 0 & \dots & 0 & - a _ {s - 1} \\ 0 & 1 & \dots & 0 & - a _ {s - 2} \\ . & . & \dots & . & . \\ . & . & \dots & . & . \\ 0 & . & \dots & 1 & - a _ {1} \end{array} \right).\tag{59}
\]

\(XI_{n}-\tau\) 的行列式可对 s 作归纳、按第一行展开来计算。我们得到 \(\det(\mathrm{X}I_{n}-\tau)=\mathrm{P}(\mathrm{X})\)。换言之，我们有 \(\mathrm{P}(\mathrm{X})=\mathrm{P}\mathrm{c}_{\mathrm{K}^{\prime\prime}/\mathrm{K}^{\prime}}(t;\mathrm{X})\)。特别地，\(\mathrm{Tr}_{\mathrm{K}^{\prime\prime}/\mathrm{K}^{\prime}}(t)=-a_{1}\)，且 \(\mathrm{N}_{\mathrm{K}^{\prime\prime}/\mathrm{K}^{\prime}}(t)=(-1)^{s}a_{s}\)。由传递公式（III, §9, No.~4, 第 548 页，推论），我们有

\[
\begin{array}{r} \mathrm{Tr} _ {\mathrm{K} ^ {\prime \prime} / \mathrm{K}} (t) = - \mathrm{Tr} _ {\mathrm{K} ^ {\prime} / \mathrm{K}} (a _ {1}), \quad \mathrm{N} _ {\mathrm{K} ^ {\prime \prime} / \mathrm{K}} (t) = (- 1) ^ {s d} \mathrm{N} _ {\mathrm{K} ^ {\prime} / \mathrm{K}} (a _ {s}), \\ \mathrm{Q} (\mathrm{X}) = \mathrm{Pc} _ {\mathrm{K} ^ {\prime \prime} / \mathrm{K}} (t; \mathrm{X}). \end{array}
\]

另一方面，\([K'' : K] = [K'' : K'][K' : K] = sd\)，因此 \(Q(X)\) 是次数为 sd 的首一多项式。由此便得引理 6。

我们来证明命题 8。由于环 B 是单的，其中心 L 是一个体（VIII, 第 121 页，推论 1）。由 VIII, 第 259 页，定理 5，B 在 A 中的交换子 \(B'\) 是一个以 L 为中心的单环，且我们有等式 \([A:K]=[B:K][B':K]\)。我们用 r 表示 \(B'\) 在 L 上的约化次数，用 s 表示 B 在 L 上的约化次数，用 d 表示 L 在 K 上的次数。我们有

\[
[ \mathrm{A}: \mathrm{K} ] = n ^ {2}, [ \mathrm{B} ^ {\prime}: \mathrm{K} ] = r ^ {2} d, [ \mathrm{B}: \mathrm{K} ] = s ^ {2} d,
\]

因而 \(n^{2}=r^{2}s^{2}d^{2}\)，即 n=rsd。

设 \(b\) 是 B 的一个元素，\(\mathrm{P}(\mathrm{X})\) 是它在 L-代数 B 上的约化特征多项式；它是次数为 \(s\) 的首一多项式。由引理 6，多项式 \(\mathrm{Q} = \mathrm{N}_{\mathrm{L}[\mathrm{X}] / \mathrm{K}[\mathrm{X}]}(\mathrm{P})\) 是次数为 \(sd\) 的首一多项式。因此多项式 \(\mathrm{R} = \mathrm{Q}^r\) 是次数为 \(rsd = n\) 的首一多项式。仍由引理 6，\(\mathrm{X}^{n - 1}\) 在 \(\mathrm{R}(\mathrm{X})\) 中的系数等于 \(-r\operatorname{Tr}_{\mathrm{L} / \mathrm{K}}(\operatorname{Trd}_{\mathrm{B} / \mathrm{L}}(b))\)，而 \(\mathrm{R}(\mathrm{X})\) 的常数项是 \((\mathrm{N}_{\mathrm{L} / \mathrm{K}}((-1)^s\mathrm{Nrd}_{\mathrm{B} / \mathrm{L}}(b)))^r = (-1)^n\mathrm{N}_{\mathrm{L} / \mathrm{K}}(\mathrm{Nrd}_{\mathrm{B} / \mathrm{L}}(b))^r\)。

由于 \([A:K]=r^{2}d[B:K]\)，左 B-模 A 是秩为 \(r^{2}d\) 的自由模（VIII, 第 124 页，命题 5）。因此我们有

\[
\mathrm{Pc} _ {\mathrm{A/K}} (b; \mathrm{X}) = \mathrm{Pc} _ {\mathrm{B/K}} (b; \mathrm{X}) ^ {d r ^ {2}}.\tag{60}
\]

由 III, §9, No.~4, 第 548 页的推论，我们有

\[
\mathrm{Pc} _ {\mathrm{B/K}} (b; \mathrm{X}) = \mathrm{N} _ {\mathrm{L[X] / K[X]}} (\mathrm{Pc} _ {\mathrm{B/L}} (b; \mathrm{X})),\tag{61}
\]

而由于 \(P(X)\) 是 b 在 L-代数 B 上的约化特征多项式，我们有

\[
\mathrm{Pc} _ {\mathrm{B/L}} (b; \mathrm{X}) = \mathrm{P(X)} ^ {s}.\tag{62}
\]

最后，由公式 (60)---(62) 以及 R(X) 的定义，我们有

\[
\mathrm{Pc} _ {\mathrm{A/K}} (b; \mathrm{X}) = \mathrm{N} _ {\mathrm{L[X]/K[X]}} (\mathrm{P(X)}) ^ {d r ^ {2} s} = \mathrm{Q(X)} ^ {d r ^ {2} s} = \mathrm{R(X)} ^ {r s d} = \mathrm{R(X)} ^ {n}\tag{63}
\]

于是 \(\mathrm{R}(\mathrm{X})\) 是 b 在 K-代数 A 上的约化特征多项式。

我们已经证明了公式 (56)。公式 (57) 与 (58) 立即由公式 (56) 和引理 6 得出，因为 \(\mathrm{X}^{n - 1}\) 在 \(\mathrm{Pcrd}_{\mathrm{A / K}}(b;\mathrm{X})\) 中的系数等于 \(-\mathrm{Trd}_{\mathrm{A / K}}(b)\)，而常数项是 \((-1)^n\mathrm{Nrd}_{\mathrm{A / K}}(b)\)。

